\section{Conclusion}
\label{sec:conclusion}

Le programme PT zêta livre un modèle spectral canonique des zéros non
triviaux de la fonction zêta de Riemann à partir d'une seule structure
interne~: la dynamique de persistance avec son axiome $\sper = 1/2$. Le
bilan final tient en cinq points.

\begin{enumerate}[leftmargin=*]
\item \textbf{L'objet opératoriel canonique est unique.} La trace
régularisée
\[
\Treg(s) \;=\; \Tr_{\mathrm{reg}}\!\big[(s - i\HPTBK)^{-1}\,\DRop(s)\big]
\;=\; \sum_{p \in \Primes_\gamma} \kappap(s)\,(\log p)\,\Big[\frac{1}{p^s - 1} + \Ap\,p^{-s}\Big]
\]
est construite à partir des amplitudes $\kappap$ de la bifurcation
$\qplus/\qminus$ et de l'opérateur de dilatation Berry--Keating PT.
Aucun paramètre libre. Identité analytique exacte dans $\Reop(s) > 1$ à
$10^{-26}$, pôles distributionnels sur $\Reop(s) = 1/2$ aux zéros de
$\Rres(s) \equiv \zeta(s) / (\zetaplus\,\zetaminus)$.

\item \textbf{La ligne critique $\Reop(s) = 1/2$ admet quatre lectures
équivalentes}~: cellule de Planck symplectique
$u_\star\,p_\star = 2\pi$, Jacobien de Haar multiplicatif,
transformation unitaire $e^{su/2}$, condition aux bords antipériodique
forcée par $T_3$. Les trois rôles du $\sper = 1/2$ ($\Tone$
arithmétique, $\Tfive$ géométrique, Haar spectral) sont
opératoriellement identifiés.

\item \textbf{Le résidu cohomologique $1/8$ est triplement déterminé.}
\[
\dfrac{1}{8} \;=\; \dfrac{\sper^{2}}{2} \;=\; \lambdaH
\;=\; \dfrac{c_1}{N_{\mathrm{corners}}} \;=\; \dfrac{\sper}{4}.
\]
La cohérence $\sper^{2}/2 = \sper/4$ admet $\sper = 0$ ou $\sper = 1/2$ ;
$\Tone$ exclut $\sper = 0$, ce qui sur-détermine $\sper = 1/2$~: c'est
une signature structurelle de $\Tone$ par la zone HP-PT, non un
substitut à $\Tone$.

\item \textbf{Dualité Higgs $\leftrightarrow \zeta$ structurelle.} Le
bifurcateur canonique $B$ (projecteur spinoriel $\qplus/\qminus$ du
canal $p = 2$ au point fixe $\mustar = 15$) produit simultanément
$m_H = v \cdot \sper$ et
$\Delta_N = \lambdaH = \sper^{2}/2$ via la chaîne
$\Tone$ + $\Ttwo$ + Haar + Weyl. Aucun autre couplage PT n'égale
$1/8$~: la coïncidence numérique pure est rejetée.

\item \textbf{Discipline de falsifiabilité.} La prédiction Dirichlet
$\Delta_N(L \bmod q) = 1/(8q)$, structurellement induite par K4, a été
testée sur les zéros LMFDB de quatre caractères quadratiques primitifs
et falsifiée ($\chi^{2}(\mathrm{PT}) = 7{,}49$ contre
$\chi^{2}(H_0) = 0{,}023$ sur 4~d.d.l.). Le résidu $1/8$ est une
identité asymptotique entre formes lisses, pas un décalage mesurable au
comptage des zéros. Le mécanisme PT décrit ici est spécifique à
$\zeta$.
\end{enumerate}

\paragraph{Position canonique.}
On a \textbf{compris} la structure des zéros de $\zeta$ ; on n'a
\textbf{pas prouvé} leur localisation stricte sur la ligne critique. Le
verrou ultime ($G_{\mathrm{HP3.a}}$~: prolongement analytique strict de
$\Rres$) est équivalent à RH au sens classique et ne semble pas plus
facile que cette dernière.

L'enrichissement principal sur Berry--Keating--Connes est la lecture
intrinsèque des coïncidences~: le $2\pi$ est le commutateur canonique,
le $1/8$ est l'auto-couplage scalaire du Higgs PT, et le spin
$\sper = 1/2$ de la ligne critique est sur-déterminé par la cohérence
cohomologique de la construction. \textbf{Aucun ingrédient n'est ad hoc.}

